% Translated from sections/04-protocol.tex on 2026-09-23. The English file is
% canonical. Change it first, then carry the change here.
\section{Protocole d'évaluation}
\label{sec:protocol}

E1, E3 et E4 n'ont pas été exécutées. Ce qui suit est le protocole sous lequel
elles le seront, fixé avant que les chiffres n'existent. Ceux de E2 et de E5,
qui ont été exécutées, sont signalés là où ils diffèrent.

\textbf{Le jeu de référence.} \goldPages{}~pages portent une transcription
corrigée. \goldPagesTresor{} viennent du \emph{Trésor}~T.1 et
\goldPagesAmadis{} du livre~13 d'\emph{Amadis de Gaule}, soit
\goldLines{}~lignes et \goldChars{}~caractères. Elles n'ont pas été annotées
pour ce rapport. Elles ont été corrigées dans Transkribus au fil de l'édition.
\file{eval/build\_gold\_set.py} les choisit parmi \goldCandidates{}~candidates.
Une page n'est retenue que si une personne l'a marquée comme terminée et
qu'elle est absente du découpage d'entraînement. Cette règle écarte
\goldRefusedTrained{}~pages. \goldRefusedUncorrected{} ne portent aucune
correction, et \goldRefusedExcluded{} relève des pièces liminaires.
\file{data/gold/MANIFEST.csv} donne une ligne par candidate avec son motif, si
bien que le jeu est vérifiable au lieu d'être affirmé. La convention de
transcription est \file{data/gold/GUIDELINES.md}, versionnée en premier.

Ce jeu remplace les vingt pages que cette section prévoyait d'abord, tirées
avec graine par strate et annotées à l'aveugle. Dix heures d'annotation étaient
la contrainte décisive, et \goldPages{}~pages déjà corrigées n'en coûtent
aucune. Quatre choses ont été cédées en échange. La référence est
\emph{postéditée et non aveugle}. Chaque page a été corrigée sur la sortie d'un
modèle, ce que le manifeste consigne page par page. Une erreur plausible y
survit donc, et le jeu penche vers le modèle affiné, le système même que E1
doit éprouver. Sa segmentation est \emph{celle de Transkribus}, ce qui rend
cette comparaison pire encore que la comparaison confondue décrite plus bas.
Son champ se limite à \emph{une œuvre et un exemplaire}, si bien que les
strates de famille et de provenance se réduisent à une seule case. Enfin,
\emph{aucune page n'a été transcrite deux fois}. Il n'y a donc pas de plancher
de bruit, et \textbf{aucune différence n'est encore assez petite pour être
écartée.} Trois pages transcrites à l'aveugle selon la même convention
fixeraient ce plancher et mesureraient le biais du même coup.

\textbf{Métriques.} Taux d'erreur caractère et taux d'erreur mot, définis dans
le code plutôt qu'en prose. L'espace compte comme un caractère, les suites
d'espaces sont réduites à une seule avant comparaison, et l'agrégation est
pondérée par caractère sur l'ensemble des pages. Chaque chiffre est donné deux
fois, brut et replié. Le repliement applique NFC, remplace le \qtr{ſ} long par
\emph{s}, ignore la casse et confond \emph{u}/\emph{v} et \emph{i}/\emph{j}. Il
est déclaré une seule fois et appliqué à l'identique à l'hypothèse et à la
référence. Il est volontairement plus étroit que le normaliseur de
l'apparieur, car tout ce qu'il effacerait en plus d'une convention de
transcription masquerait une vraie erreur de reconnaissance. L'incertitude est
estimée par bootstrap par percentiles sur les pages, avec 10\,000~tirages, à
95\,\%~\cite{efron1994bootstrap}. Sur les pages et non sur les caractères, car
une page mal encrée ou mal segmentée échoue en bloc. Rééchantillonner les
caractères donnerait un intervalle plus étroit que ce que les données
permettent.

\textbf{Systèmes de comparaison.} E1 opposera le modèle affiné à CATMuS-Print
sans affinage, à un modèle publié choisi pour la première modernité et à
Transkribus. La segmentation sera la même pour les trois premiers. Elle ne peut
pas l'être pour Transkribus, qui segmente de son côté. Cette comparaison mêle
donc reconnaissance et mise en page, et elle est rapportée comme confondue.

\textbf{Deux préenregistrements.} Le critère de E5 a été écrit le 15~septembre
2026, avant l'évaluation de l'apparieur, et il est cité mot pour mot dans
\file{docs/pre-registration/2026-09-15-e5-alignment-gate.md}. Le protocole de
E2 a été enregistré le 22~septembre 2026, avant son exécution. Il fixe la base
de sondage, l'échantillon et sa graine, le traitement, et la liste figée des
chiffres à rapporter. E2 est la seule expérience de ce rapport qui pouvait être
relancée à volonté. C'est précisément pourquoi ses choix sont datés avant ses
chiffres.

\textbf{Ce que le protocole ne peut pas faire.} Mesurer la passe de correction
contre l'absence de correction, en variation du taux d'erreur, exige une
transcription de référence pour les pages que E2 a tirées. Le jeu de référence
ne les couvre pas. Il porte sur le \emph{Trésor} et sur un extrait
d'\emph{Amadis}, alors que la base de sondage de E2 est la sortie de
reconnaissance propre aux \ocrBooks{}~livres. E2 rapporte donc ce qu'a
fait la règle d'acceptation, sans prétendre que les modifications retenues
sont des améliorations. Cette limite a été enregistrée à l'avance, et non
découverte après coup.
